% =============================================================================
% 07-process-operations-notes.tex — Lecture Notes for Lecture 7:
% Operations on Processes — Creation, Termination, fork/exec
%
% DERIVED ARTIFACT. Source of truth: Slides/CS502/07-process-operations.tex.
% Generated by /lecture-notes; do not hand-edit content independently of the
% Beamer source without re-running /qa-notes afterward (single-source-of-truth.md).
% Structured as a textbook chapter (numbered sections/figures/examples,
% topic-organized rather than a frame-by-frame transcript).
%
% Anchor reading: Silberschatz, Galvin & Gagne, Operating System Concepts,
%                 8th ed. — Ch. 3 §3.3 (pp. 110-116); Stallings 6e Ch. 3 §3.4
%                 (p. 135); Bovet & Cesati 3e Ch. 3 (pp. 79-130) and Ch. 20
%                 (pp. 808-834). All page-verified via the indexes under
%                 master_supporting_docs/CS502/supporting_books/.
%
% Every fork() output count below was verified by compiling and running the
% corresponding C program (gcc) on 2026-10-29.
%
% Compile with: cd Notes/CS502 &&
%   TEXINPUTS=../../Preambles:$TEXINPUTS xelatex -interaction=nonstopmode 07-process-operations-notes.tex  (x3)
% =============================================================================

\documentclass[11pt]{article}
\usepackage[margin=1in]{geometry}
\usepackage{array}
\input{../../Preambles/header}
\coursecode{CS502}

\renewcommand{\thesection}{7.\arabic{section}}
\renewcommand{\thefigure}{7.\arabic{figure}}
\renewcommand{\thetable}{7.\arabic{table}}

\newtheorem{example}{Example}
\renewcommand{\theexample}{7.\arabic{example}}

\newenvironment{definitionbox}[1]{%
  \par\noindent\textbf{Definition (#1).}\ \itshape
}{\par\vspace{0.3em}}

\setlength{\emergencystretch}{2.5em}

\title{Operations on Processes\\\large Lecture Notes --- Lecture 7\\
  \normalsize Creation, Termination, \texttt{fork}/\texttt{exec}}
\author{Auqib Hamid Lone\\[0.3em]
  \emph{Department of Computer Science \& Engineering}, \emph{GCET Ganderbal}}
\date{\today}

\begin{document}
\maketitle

\section{Who Actually Runs a Command?}

Type \texttt{ls} at a shell prompt. A directory listing appears and the prompt
returns. Two facts about this ordinary event are worth stating carefully. The
shell did not turn into \texttt{ls} --- it is still there, waiting for the next
command. And \texttt{ls} did not exist as a process a moment earlier. So a new
process came into being, ran, and disappeared, and the shell knew exactly when
it had finished. This chapter explains how.

The short answer is that every process except the very first is created by
another process. The shell creates a \key{child}, the child turns itself into
\texttt{ls}, and the shell \key{waits} for it. Three system calls do the work
--- \texttt{fork}, \texttt{exec}, and \texttt{wait} --- and a fourth,
\texttt{exit}, ends the child. We follow one shell command from creation to
termination as the running thread of the chapter.

The chapter answers four questions in order: how a process is created (parents,
children, and the process tree); what \texttt{fork} returns and how many
processes a program produces (counting forks); how a child becomes a different
program (the \texttt{exec} family); and how a process ends and who cleans up
after it (\texttt{exit}, \texttt{wait}, zombies, and orphans). Lecture~5 gave us
the PCB and Lecture~6 moved a process between states; here we create and
destroy the PCB itself.

\medskip\noindent\textit{Reading:} Silberschatz et al., Ch.~3 \S3.3,
pp.~110--116 \cite{SilberschatzOSC_operating_system_concepts}; Stallings, Ch.~3
\S3.4, p.~135 \cite{StallingsOS_internals_design_principles}.

\section{Process Creation}

\subsection{Parents, children, and the process tree}

A process is not conjured by the kernel out of nothing: it is copied from a
parent that already exists. When one process creates another, the creator is
the parent and the new process is the child; the child may in turn create
children of its own. Recording every creation as an arrow from parent to child
produces a tree.

\begin{definitionbox}{Parent, child, process tree}
A \key{parent process} is a process that issues a creation call; the new
process is its \key{child}. Repeated creation forms a \key{process tree} rooted
at the first process.
\end{definitionbox}

Each process is identified by a unique integer, its \key{process identifier}
(pid). Figure~\ref{fig:tree} shows a typical tree on a UNIX-like system: the
first process, \texttt{init} with pid 1, has started a login program, an SSH
daemon, and a scheduler daemon; each login path has started a shell; and one
shell has started \texttt{ls} and \texttt{vim}. The pids are illustrative.

\begin{figure}[h]
\centering
% Coordinate map (P2): level y=0: init (0,0)
  %   level y=-1.5: login (-4,-1.5), sshd (0,-1.5), cron (4,-1.5)
  %   level y=-3.0: bash (-4,-3.0), bash (0,-3.0)
  %   level y=-4.5: ls (-5.6,-4.5), vim (-2.4,-4.5)
  % Boxes 2.2cm x 0.7cm; nearest horizontal gap 3.2-2.2 = 1.0cm (clear).
  % Edges run box-bottom to box-top only; no edge crosses a box.
  \begin{tikzpicture}[every node/.style={font=\small},
      pnode/.style={dag-node, minimum width=2.2cm, minimum height=0.7cm}]
    \node[pnode] (init) at (0,0) {\texttt{init} (1)};
    \node[pnode] (login) at (-4,-1.5) {\texttt{login} (812)};
    \node[pnode] (sshd) at (0,-1.5) {\texttt{sshd} (640)};
    \node[pnode] (cron) at (4,-1.5) {\texttt{cron} (701)};
    \node[pnode] (bash1) at (-4,-3.0) {\texttt{bash} (905)};
    \node[pnode] (bash2) at (0,-3.0) {\texttt{bash} (1220)};
    \node[pnode] (ls) at (-5.6,-4.5) {\texttt{ls} (1310)};
    \node[pnode] (vim) at (-2.4,-4.5) {\texttt{vim} (1288)};
    \draw[observed-edge] (init.south) -- (login.north);
    \draw[observed-edge] (init.south) -- (sshd.north);
    \draw[observed-edge] (init.south) -- (cron.north);
    \draw[observed-edge] (login.south) -- (bash1.north);
    \draw[observed-edge] (sshd.south) -- (bash2.north);
    \draw[observed-edge] (bash1.south) -- (ls.north);
    \draw[observed-edge] (bash1.south) -- (vim.north);
  \end{tikzpicture}
\caption{A process tree. Each arrow is one creation; pids are illustrative ---
cf.\ Silberschatz et al., Ch.~3 \S3.3, pp.~110--116
\cite{SilberschatzOSC_operating_system_concepts}.}
\label{fig:tree}
\end{figure}

\subsection{What a parent must decide}

Creating a child forces three design decisions, and every operating system
answers them in some way \cite{SilberschatzOSC_operating_system_concepts}:

\begin{enumerate}
  \item \textbf{Resources.} The child may share all of the parent's resources
    (open files, memory), a subset of them, or none.
  \item \textbf{Execution.} The parent may continue to run concurrently with
    the child, or it may wait until some or all of its children terminate.
  \item \textbf{Address space.} The child may be a duplicate of the parent
    (the same program and data), or it may have a new program loaded into it.
\end{enumerate}

\noindent UNIX answers the address-space question in two separate steps:
\texttt{fork} duplicates the parent, and \texttt{exec} loads a new program.

\subsection{\texttt{fork}: one call, two returns}

\begin{definitionbox}{\texttt{fork}}
\texttt{fork()} creates a new process whose address space is a copy of the
caller's. Both processes continue at the instruction after the call.
\end{definitionbox}

Because both processes resume at the same point in the same code, they need a
way to tell themselves apart. The only difference between them at that moment
is the value \texttt{fork} returns, summarized in Table~\ref{tab:fork-return}.

\begin{table}[h]
\centering
\renewcommand{\arraystretch}{1.3}
\begin{tabular}{@{}p{0.32\linewidth}p{0.52\linewidth}@{}}
\toprule
\textbf{Who sees it} & \textbf{Return value of \texttt{fork()}} \\
\midrule
the child & \texttt{0} \\
the parent & the child's pid (a positive integer) \\
the parent, on failure & \texttt{-1}, and no child is created \\
\bottomrule
\end{tabular}
\caption{The three possible return values of \texttt{fork}.}
\label{tab:fork-return}
\end{table}

\noindent The return value is the \emph{only} way the two copies can
distinguish themselves. A child that needs its own pid calls \texttt{getpid()};
it learns its parent's pid with \texttt{getppid()}.

\begin{example}[The canonical \texttt{fork}/\texttt{exec}/\texttt{wait}]
\label{ex:canonical}
The following fragment has the same structure as the standard textbook example
\cite{SilberschatzOSC_operating_system_concepts}:
\begin{verbatim}
pid_t pid = fork();
if (pid < 0) {            /* error  */
    perror("fork");
    exit(1);
} else if (pid == 0) {    /* child  */
    execlp("ls", "ls", NULL);
    perror("exec");  /* on failure */
    exit(1);
} else {                  /* parent */
    wait(NULL);
    printf("Child complete\n");
}
\end{verbatim}
Trace it step by step.
\begin{enumerate}
  \item \texttt{fork} executes once and returns twice: once in the parent with
    the child's pid, once in the child with \texttt{0}.
  \item In the child, \texttt{pid == 0}, so the child takes the middle branch
    and calls \texttt{execlp}, replacing itself with the \texttt{ls} program.
    The \texttt{perror} and \texttt{exit(1)} that follow run only if the
    \texttt{exec} fails.
  \item In the parent, \texttt{pid} is positive, so the parent takes the
    \texttt{else} branch and blocks in \texttt{wait}.
  \item When \texttt{ls} finishes, \texttt{wait} returns and the parent prints
    \texttt{Child complete}.
\end{enumerate}
\end{example}

Figure~\ref{fig:timeline} draws the same program as a timeline. The parent's
line is interrupted by the \texttt{wait} box: the parent is blocked there while
the child runs. The child's line begins at the \texttt{fork} point, passes
through \texttt{exec("ls")}, and ends at \texttt{exit()}; the arrow back up to
the parent carries the exit status and wakes the parent.

\begin{figure}[h]
\centering
% Coordinate map (P2):
  %   parent line y=0 from x=-6.5 to x=6.5; child line y=-1.8 from x=-3.5 to x=3.5
  %   fork point (-3.5,0) -> child start (-3.5,-1.8)
  %   exec box centered (-1.2,-1.8), 2.2cm x 0.6cm
  %   exit point (3.5,-1.8) -> resume point (3.5,0)
  %   wait box centered (0.5,0), 3.0cm x 0.6cm (spans -1.0..2.0)
  %   labels: "fork()" above (-3.5,0.45); "exit()" right of (3.5,-1.8) at x=4.4;
  %           "parent" left at x=-6.9 y=0; "child" left at x=-4.2 y=-1.8
  % Arrow segments start/end at box edges; no label sits on a line.
  \begin{tikzpicture}[scale=0.92, transform shape, every node/.style={font=\small}]
    \draw[observed-edge] (-6.5,0) -- (-1.5,0);
    \node[dag-node, minimum width=3.0cm, minimum height=0.6cm] (wait) at (0.5,0) {\texttt{wait()}};
    \draw[observed-edge] (2.0,0) -- (6.5,0);
    \node[anchor=east] at (-6.6,0) {parent};
    \node at (-3.5,0.45) {\texttt{fork()}};
    \fill[primary-blue] (-3.5,0) circle (2.5pt);
    \draw[observed-edge] (-3.5,0) -- (-3.5,-1.65);
    \draw[thick, draw=jet] (-3.5,-1.8) -- (-2.3,-1.8);
    \node[dag-node, minimum width=2.2cm, minimum height=0.6cm] (exec) at (-1.2,-1.8) {\texttt{exec("ls")}};
    \draw[thick, draw=jet] (-0.1,-1.8) -- (3.5,-1.8);
    \node[anchor=east] at (-3.7,-1.8) {child};
    \node[anchor=west] at (3.65,-1.8) {\texttt{exit()}};
    \draw[observed-edge] (3.5,-1.8) -- (3.5,-0.15);
    \node[anchor=west, text=neutral] at (3.6,-0.9) {status};
  \end{tikzpicture}
\caption{\texttt{fork}, \texttt{exec}, \texttt{exit}, and \texttt{wait} on a
timeline. The parent is blocked in \texttt{wait} while the child runs; the
child's \texttt{exit} wakes it and delivers a status.}
\label{fig:timeline}
\end{figure}

\subsection{Two copies, two variables}

After \texttt{fork}, the parent and child have \emph{separate} address spaces
whose contents are initially identical. A write in one is invisible to the
other.

\begin{example}[Separate copies of a variable]
\label{ex:copies}
Consider \texttt{int x = 10; pid = fork();} followed by:
\begin{itemize}
  \item child: \texttt{x += 5; printf("child x=\%d", x);}
  \item parent: \texttt{wait(NULL); x -= 5; printf("parent x=\%d", x);}
\end{itemize}
Both processes start with their own \texttt{x} equal to 10. The child adds 5
to its copy and prints \texttt{child x=15}. The parent waits for the child,
then subtracts 5 from \emph{its} copy, which is still 10, and prints
\texttt{parent x=5}. (This output was verified by running the program.)
\end{example}

The variable has the same name and the same virtual address in both
processes, but the two names refer to different physical copies. What
\emph{is} shared after \texttt{fork} is the set of open file descriptors: the
child's descriptors point to the same open-file entries as the parent's.

A consequence follows immediately. Suppose a child wants to report the answer
to a computation by storing it in a global variable that the parent reads
after \texttt{wait}. That cannot work: the global lives in the child's copy of
the address space, and the parent's copy still holds the old value. The child
can return a small integer through its \texttt{exit} status; anything larger
requires inter-process communication, the subject of Lectures~14--15.

\section{Counting Forks}

Every \texttt{fork} doubles the number of processes that execute the next
line. Counting the processes a program creates is therefore a matter of
drawing a tree, not of guessing.

\subsection{Sequential forks}

\begin{example}[Three forks in a row]
\label{ex:three}
Count the lines printed by
\texttt{fork(); fork(); fork(); printf("hi\textbackslash n");}.
\begin{enumerate}
  \item Before any fork there is one process, P.
  \item The first \texttt{fork} creates C1. Now 2 processes execute the
    second \texttt{fork}.
  \item Each of the 2 forks: P creates C2, C1 creates C3. Now 4 processes
    execute the third \texttt{fork}.
  \item Each of the 4 forks: P creates C4, C2 creates C5, C1 creates C6, C3
    creates C7. Now 8 processes execute the \texttt{printf}.
\end{enumerate}
So \key{8} lines are printed (Figure~\ref{fig:forktree}). In general, $n$
sequential forks give $2^n$ processes, of which $2^n - 1$ are new children.
Running the program confirms 8 lines for $n=3$ and 16 for $n=4$.
\end{example}

\begin{figure}[h]
\centering
% Coordinate map (P2): binary tree, levels at y=0, -1.1, -2.2, -3.3
  %   level 0: P (0,0)
  %   level 1: x = -3.2, 3.2
  %   level 2: x = -4.8, -1.6, 1.6, 4.8
  %   level 3: x = -5.6,-4.0,-2.4,-0.8,0.8,2.4,4.0,5.6 (circles r=0.32cm)
  % Nearest leaf gap 1.6-0.64 = 0.96cm. Level labels at x=7.2, clear of x=5.92.
  \begin{tikzpicture}[scale=0.9, transform shape, every node/.style={font=\footnotesize},
      proc/.style={circle, draw=primary-blue, fill=light-bg, minimum size=0.64cm, inner sep=0pt}]
    \node[proc] (a) at (0,0) {P};
    \node[proc] (b1) at (-3.2,-1.1) {P};
    \node[proc] (b2) at (3.2,-1.1) {C1};
    \foreach \x/\n/\l in {-4.8/c1/P, -1.6/c2/C2, 1.6/c3/C1, 4.8/c4/C3}
      \node[proc] (\n) at (\x,-2.2) {\l};
    \foreach \x/\n/\l in {-5.6/d1/P, -4.0/d2/C4, -2.4/d3/C2, -0.8/d4/C5,
                          0.8/d5/C1, 2.4/d6/C6, 4.0/d7/C3, 5.6/d8/C7}
      \node[proc] (\n) at (\x,-3.3) {\l};
    \draw[observed-edge] (a) -- (b1); \draw[observed-edge] (a) -- (b2);
    \draw[observed-edge] (b1) -- (c1); \draw[observed-edge] (b1) -- (c2);
    \draw[observed-edge] (b2) -- (c3); \draw[observed-edge] (b2) -- (c4);
    \draw[observed-edge] (c1) -- (d1); \draw[observed-edge] (c1) -- (d2);
    \draw[observed-edge] (c2) -- (d3); \draw[observed-edge] (c2) -- (d4);
    \draw[observed-edge] (c3) -- (d5); \draw[observed-edge] (c3) -- (d6);
    \draw[observed-edge] (c4) -- (d7); \draw[observed-edge] (c4) -- (d8);
    \node[anchor=west, text=neutral] at (6.4,-1.1) {after 1st: 2};
    \node[anchor=west, text=neutral] at (6.4,-2.2) {after 2nd: 4};
    \node[anchor=west, text=neutral] at (6.4,-3.3) {after 3rd: 8};
  \end{tikzpicture}
\caption{The process tree for three sequential \texttt{fork} calls. Each level
doubles the number of processes: 2, then 4, then 8.}
\label{fig:forktree}
\end{figure}

\begin{example}[\texttt{printf} inside the loop]
\label{ex:loop}
Count the lines printed by
\texttt{for (i = 0; i < 3; i++) \{ fork(); printf("A\textbackslash n"); \}}.
Here the print is \emph{inside} the loop, so we count per iteration.
\begin{enumerate}
  \item Iteration $i=0$: 1 process forks into 2; each prints once: 2 lines.
  \item Iteration $i=1$: 2 processes fork into 4; each prints: 4 lines.
  \item Iteration $i=2$: 4 processes fork into 8; each prints: 8 lines.
\end{enumerate}
The total is $2 + 4 + 8 = 14$ lines of \texttt{A}, from 8 processes in total.
(Verified by running it.)
\end{example}

\subsection{Short-circuit operators prune the tree}

In C, \texttt{a \&\& b} evaluates \texttt{b} only if \texttt{a} is non-zero,
and \texttt{a || b} evaluates \texttt{b} only if \texttt{a} is zero. Since
\texttt{fork()} is non-zero in the parent and zero in the child, these
operators decide \emph{which copy} goes on to fork again.

\begin{example}[\texttt{(fork() \&\& fork()) || fork()}]
\label{ex:shortcircuit}
Count the lines printed by
\texttt{(fork() \&\& fork()) || fork(); printf("y\textbackslash n");}.
\begin{enumerate}
  \item First \texttt{fork}: the parent P sees a non-zero value; the child C1
    sees zero.
  \item P evaluates the second \texttt{fork}: P sees non-zero, the new child
    C2 sees zero. P's \texttt{\&\&} is true, so P skips the \texttt{||}.
  \item C1's \texttt{\&\&} is already false (its first operand was zero), so
    C1 evaluates the third \texttt{fork}, creating C3.
  \item C2's \texttt{\&\&} is false (its second operand was zero), so C2 also
    evaluates the third \texttt{fork}, creating C4.
\end{enumerate}
C3 and C4 see zero from the third \texttt{fork} but have nothing left to
evaluate. The processes are P, C1, C2, C3, C4, so \key{5} lines are printed.
(Verified.)
\end{example}

\begin{example}[\texttt{if (fork() \&\& fork()) fork()}]
\label{ex:iffork}
Count the lines printed by
\texttt{if (fork() \&\& fork()) fork(); printf("x\textbackslash n");}.
\begin{enumerate}
  \item First \texttt{fork}: P and C1. C1 sees 0, so its condition is false and
    it skips the body.
  \item P runs the second \texttt{fork}: P and C2. C2 sees 0 and skips the body.
  \item P's condition is true, so P runs the body's \texttt{fork}, creating C3.
\end{enumerate}
The processes are P, C1, C2, C3: \key{4} lines. (Verified.)
\end{example}

\subsection{The buffer that gets copied}

Lecture~3 showed that \texttt{printf} writes into a user-space buffer and that
the \texttt{write} system call comes later. \texttt{fork} copies the
\emph{whole} address space --- including any bytes still sitting in that
buffer.

\begin{example}[Duplicated output]
\label{ex:buffer}
The program
\begin{center}\texttt{printf("hello "); fork(); printf("world\textbackslash n");}\end{center}
\noindent prints \texttt{hello world} \emph{twice}. The string \texttt{"hello "} is still
unflushed in the stdio buffer when \texttt{fork} runs, so it is duplicated into
the child; each process then appends \texttt{"world\textbackslash n"} and
flushes its own buffer. (Verified.) The fix is to call \texttt{fflush(stdout)}
before \texttt{fork}. The lesson is that \texttt{fork} copies \emph{state}, not
just code.
\end{example}

\section{The \texttt{exec} Family}

A copy of the shell is useless until it stops being the shell. \texttt{exec}
throws the old program away and keeps the process.

\begin{definitionbox}{\texttt{exec}}
An \texttt{exec} call replaces the calling process's memory image --- code,
data, heap, stack --- with a new program loaded from an executable file, and
starts it at its entry point.
\end{definitionbox}

Several things survive the replacement. The \key{pid stays the same}, and so
does the parent relationship. Open file descriptors stay open unless they are
marked close-on-exec. What does not survive is the old code itself, which is
why a successful \texttt{exec} \key{never returns}: the code that called it no
longer exists. On failure \texttt{exec} returns \texttt{-1} and the old program
continues. The Linux implementation is \texttt{execve}
(Bovet \& Cesati, Ch.~20, pp.~808--834 \cite{BovetCesati_understanding_linux_kernel}).

\subsection{Reading the \texttt{exec} names}

The family's names are built from letters, each with a fixed meaning
(Table~\ref{tab:exec}).

\begin{table}[h]
\centering
\renewcommand{\arraystretch}{1.3}
\begin{tabular}{@{}p{0.12\linewidth}p{0.76\linewidth}@{}}
\toprule
\textbf{Letter} & \textbf{Meaning} \\
\midrule
\texttt{l} & arguments as a \emph{list}: \texttt{execl("/bin/ls", "ls", "-l", NULL)} \\
\texttt{v} & arguments as a \emph{vector} (array): \texttt{execv(path, argv)} \\
\texttt{p} & search the \texttt{PATH} for the file: \texttt{execlp("ls", ...)} \\
\texttt{e} & pass an explicit \emph{environment}: \texttt{execle}, \texttt{execve} \\
\bottomrule
\end{tabular}
\caption{Decoding the \texttt{exec} family names.}
\label{tab:exec}
\end{table}

\noindent All of these are library wrappers; on Linux the one system call
underneath is \texttt{execve}. Argument lists end with \texttt{NULL}, and by
convention \texttt{argv[0]} is the program's own name.

\begin{example}[Code after \texttt{exec}]
\label{ex:afterexec}
The child executes
\texttt{execlp("echo", "echo", "from-exec", NULL); printf("never\textbackslash n");}
while the parent executes \texttt{wait(NULL); printf("parent\textbackslash n");}.
\begin{enumerate}
  \item The child's image is replaced by \texttt{echo}, which prints
    \texttt{from-exec} and exits.
  \item The \texttt{printf("never")} line belonged to the old image. It is
    gone, so it never runs.
  \item The parent's \texttt{wait} returns, and the parent prints
    \texttt{parent}.
\end{enumerate}
The output is \texttt{from-exec} followed by \texttt{parent} --- two lines,
never three. (Verified.)
\end{example}

\subsection{Why split creation into \texttt{fork} and \texttt{exec}?}

Between the two calls, the child is still running \emph{the shell's code}, but
in its own process. That gap is where it can rearrange its own environment
before the new program starts.

\begin{example}[The shell running \texttt{ls > out.txt}]
\label{ex:redirect}
\begin{enumerate}
  \item The shell calls \texttt{fork}.
  \item The child opens the target:
    \texttt{fd = open("out.txt", O\_WRONLY|O\_CREAT|O\_TRUNC)}.
  \item The child calls \texttt{dup2(fd, 1)} and then \texttt{close(fd)};
    standard output (descriptor 1) now refers to the file.
  \item The child calls \texttt{execlp("ls", "ls", NULL)}. \texttt{ls} writes
    to descriptor 1 as usual and never knows about the file.
  \item The shell calls \texttt{wait}.
\end{enumerate}
\end{example}

\noindent Redirection, pipes, and changed working directories all happen in
this gap.

\subsection{Is copying the whole address space wasteful?}

A child that calls \texttt{exec} immediately discards the copy that
\texttt{fork} just made, so a literal copy would be wasted work. Modern kernels
avoid it with \key{copy-on-write}: parent and child share the pages read-only,
and a page is copied only when one of them writes to it. \texttt{fork} then
costs roughly a copy of the page tables and the PCB, not of the memory. Linux
also provides \texttt{vfork} and \texttt{clone} as variants of process creation
(Bovet \& Cesati, Ch.~3, pp.~79--130
\cite{BovetCesati_understanding_linux_kernel}). Copy-on-write is a
virtual-memory mechanism, and we return to it in Lecture~26.

\subsection{Another design: one-step creation}

UNIX is not the only answer. Table~\ref{tab:unixwin} contrasts it with the
Windows design \cite{SilberschatzOSC_operating_system_concepts}.

\begin{table}[h]
\centering
\renewcommand{\arraystretch}{1.3}
\begin{tabular}{@{}p{0.42\linewidth}p{0.46\linewidth}@{}}
\toprule
\textbf{UNIX: two steps} & \textbf{Windows: one step} \\
\midrule
\texttt{fork} duplicates, then \texttt{exec} loads &
\texttt{CreateProcess} loads a named program into a new process \\
Child customizes itself between the calls &
Customization passed as parameters to the single call \\
Parent waits with \texttt{wait} &
Parent waits with \texttt{WaitForSingleObject} \\
\bottomrule
\end{tabular}
\caption{Two-step versus one-step process creation.}
\label{tab:unixwin}
\end{table}

\noindent Both designs answer the same three parental decisions (resources,
execution, address space); they differ in \emph{where} the answers are
supplied.

\section{Process Termination}

Creation is half the life cycle. When a process ends, something must collect
what it leaves behind.

\subsection{Normal termination: \texttt{exit} and \texttt{wait}}

A process terminates when it finishes its last statement or calls
\texttt{exit(status)}. The kernel then releases its memory, open files, and
other resources, but keeps the small integer \texttt{status} for the parent.
The parent collects it with \texttt{wait}, which returns the pid of the
terminated child; \texttt{waitpid} waits for one named child, whereas
\texttt{wait} waits for any.

\begin{example}[Collecting an exit status]
\label{ex:status}
The child calls \texttt{exit(7)}. The parent calls
\texttt{waitpid(pid, \&st, 0)}. When it returns, \texttt{WIFEXITED(st)} is true
(the child exited normally) and \texttt{WEXITSTATUS(st)} is 7. (Verified.)
\end{example}

\subsection{When a parent terminates a child}

A parent may also end a child with a system call --- \texttt{kill} on UNIX,
\texttt{TerminateProcess} on Windows. Only a parent, or a suitably privileged
user, may do so. Typical reasons \cite{SilberschatzOSC_operating_system_concepts}
are that the child has exceeded its allocated resources, that the task assigned
to it is no longer required, or that the parent is exiting and the system does
not allow a child to continue without its parent.

\begin{definitionbox}{Cascading termination}
On systems that do not allow a child to outlive its parent, terminating a
parent terminates all of its descendants. This \key{cascading termination} is
initiated by the operating system.
\end{definitionbox}

\subsection{Zombies and orphans}

The gap between a child's \texttt{exit} and its parent's \texttt{wait} creates
two bookkeeping states. Their treatment here is the standard UNIX one
(Bovet \& Cesati, Ch.~3, destroying processes, pp.~79--130
\cite{BovetCesati_understanding_linux_kernel}).

\begin{definitionbox}{Zombie}
A \key{zombie} is a process that has terminated but whose parent has not yet
called \texttt{wait}. Its resources are freed, but its PCB entry (pid and exit
status) remains so that the parent can collect it.
\end{definitionbox}

\begin{definitionbox}{Orphan}
An \key{orphan} is a process whose parent terminated without waiting for it.
UNIX re-parents an orphan to \texttt{init} (pid 1), which calls \texttt{wait}
periodically.
\end{definitionbox}

Every child is briefly a zombie, between its \texttt{exit} and its parent's
\texttt{wait}; that is normal. The problem is a parent that never waits. Many
zombies exhaust the pid table even though they use no memory. A zombie is a
bookkeeping state, not a running program. Figure~\ref{fig:lifecycle} shows both
paths: the solid path is the normal case (\texttt{fork} creates, \texttt{exit}
makes a zombie, the parent's \texttt{wait} reaps it); the dashed path is the
case in which the parent exits first, the child becomes an orphan of
\texttt{init}, and \texttt{init}'s \texttt{wait} reaps it.

\begin{figure}[h]
\centering
% Coordinate map (P2): boxes 2.6cm x 0.8cm
  %   created (-6,0), running (-2,0), zombie (2,0), reaped (6,0)
  %   orphan (2,-2.2): adopted-by-init path from running
  % Horizontal gaps 4.0-2.6 = 1.4cm; edge labels placed above the arrows
  % (y=+0.65), never on the line. Orphan arrow: running.south -> orphan.west.
  % "init waits" sits below the orphan->reaped horizontal leg (y=-2.2), at y=-2.45.
  \begin{tikzpicture}[scale=0.88, transform shape, every node/.style={font=\small},
      st/.style={dag-node, minimum width=2.6cm, minimum height=0.8cm}]
    \node[st] (new) at (-6,0) {created};
    \node[st] (run) at (-2,0) {running};
    \node[st] (zom) at (2,0) {zombie};
    \node[st] (gone) at (6,0) {reaped};
    \node[st] (orph) at (2,-2.2) {orphan of \texttt{init}};
    \draw[observed-edge] (new) -- (run);
    \draw[observed-edge] (run) -- (zom);
    \draw[observed-edge] (zom) -- (gone);
    \node[text=primary-gold] at (-4,0.65) {\texttt{fork}};
    \node[text=primary-gold] at (0,0.65) {\texttt{exit}};
    \node[text=primary-gold] at (4,0.65) {parent \texttt{wait}};
    \draw[counterfactual-edge] (run.south) |- (orph.west);
    \node[text=neutral, anchor=east] at (-2.15,-1.3) {parent exits first};
    \draw[counterfactual-edge] (orph.east) -| (gone.south);
    \node[text=neutral, anchor=north] at (4.65,-2.45) {\texttt{init} waits};
  \end{tikzpicture}
\caption{The process life cycle at the boundary. Solid: the normal path.
Dashed: the parent exits first and \texttt{init} adopts and reaps the child.}
\label{fig:lifecycle}
\end{figure}

\begin{example}[A server that never waits]
\label{ex:server}
A server forks one child per request and never calls \texttt{wait}; each child
finishes in milliseconds. Every finished child becomes a zombie, so its pid and
status stay in the process table. Memory is released, but after enough requests
the pids run out. The fix is to reap the children --- call \texttt{wait} or
\texttt{waitpid} as they finish, for example in a loop with \texttt{WNOHANG}.
\end{example}

\subsection{Putting it together: the shell loop}

\begin{example}[The shell loop]
\label{ex:shell}
A shell repeats five steps.
\begin{enumerate}
  \item Print the prompt and read a command line, e.g.\ \texttt{ls -l}.
  \item Call \texttt{pid = fork()}.
  \item In the child (\texttt{pid == 0}): set up any redirection, then call
    \texttt{execvp("ls", argv)}; if \texttt{exec} returns, print an error and
    call \texttt{exit(127)}.
  \item In the parent: call \texttt{waitpid(pid, \&status, 0)} --- unless the
    line ended with \texttt{\&}, in which case return to step 1 immediately.
  \item Repeat.
\end{enumerate}
For each foreground command the process-control calls are \texttt{fork}
(shell), \texttt{execvp} (child), \texttt{exit} (child), and \texttt{wait}
(shell): \key{4} calls, all in the process-control category of Lecture~3.
\end{example}

\section{Common Misconceptions}

Table~\ref{tab:misc} pairs five common errors with their corrections.

\begin{table}[h]
\centering
\small
\renewcommand{\arraystretch}{1.3}
\begin{tabular}{@{}>{\raggedright\arraybackslash}p{0.44\linewidth}>{\raggedright\arraybackslash}p{0.48\linewidth}@{}}
\toprule
\textbf{Wrong} & \textbf{Right} \\
\midrule
``\texttt{fork} returns the child's pid to the child.'' & The child gets \texttt{0}; the parent gets the child's pid. \\
``Parent and child share variables.'' & They have separate copies; only descriptors are shared. \\
``\texttt{exec} creates a new process.'' & \texttt{exec} replaces the image; the pid is unchanged. \\
``Code after a successful \texttt{exec} runs when it finishes.'' & A successful \texttt{exec} never returns. \\
``A zombie is still using CPU and memory.'' & A zombie holds only its pid and exit status. \\
\bottomrule
\end{tabular}
\caption{Common misconceptions about process operations.}
\label{tab:misc}
\end{table}

\section{Summary}

\begin{itemize}
  \item Processes form a \key{tree}: every process but the first has a parent.
  \item \texttt{fork} returns \key{twice}: \texttt{0} in the child, the child's
    pid in the parent, \texttt{-1} on failure.
  \item $n$ sequential forks give $2^n$ processes; short-circuit operators and
    loops change the count, so draw the tree.
  \item \texttt{exec} keeps the process and replaces the program; it returns
    only on failure.
  \item \texttt{exit} delivers a status and \texttt{wait} collects it. An
    unreaped child is a \key{zombie}; a parentless one is an \key{orphan}
    adopted by \texttt{init}.
\end{itemize}

\section{Looking Ahead and Further Reading}

This chapter covered how processes are born, transformed, and reaped. The next
lecture asks what happens once many processes are ready: which one gets the
CPU? It introduces the scheduling criteria (TAT, WT, RT) and the first two
policies, FCFS and SJF. Read Silberschatz et al., Ch.~5 \S5.1--5.3
\cite{SilberschatzOSC_operating_system_concepts}, and Stallings, Ch.~9
\S9.1--9.2 \cite{StallingsOS_internals_design_principles}.

For this chapter's material:
\begin{itemize}
  \item Silberschatz et al., 8e, Ch.~3 \S3.3 (operations on processes:
    creation, termination, \texttt{fork}/\texttt{exec}, pp.~110--116)
    \cite{SilberschatzOSC_operating_system_concepts}.
  \item Stallings, 6e, Ch.~3 \S3.4 (process control, p.~135)
    \cite{StallingsOS_internals_design_principles}.
  \item Bovet \& Cesati, 3e, Ch.~3 (creating and destroying processes,
    pp.~79--130) and Ch.~20 (program execution, pp.~808--834)
    \cite{BovetCesati_understanding_linux_kernel}.
  \item Bach, \emph{The Design of the UNIX Operating System} --- chapter-level
    reference for the UNIX process view
    \cite{Bach1986_design_unix_operating_system}.
\end{itemize}
Page numbers are verified via each book's \texttt{index.md} under the course's
supporting-books directory; Bach is not indexed.

\section*{Self-Check}
\addcontentsline{toc}{section}{Self-Check}

\subsection*{Exercises}

\begin{enumerate}
  \item State the three possible return values of \texttt{fork} and who sees
        each. Why is the return value the only way the two copies can tell
        themselves apart?
  \item How many lines does
        \\\texttt{for (i = 0; i < 4; i++) fork(); printf("p\textbackslash n");}\\
        print, and how many new processes are created?
  \item How many lines does
        \\\texttt{if (fork() == 0) fork(); else \{ fork(); fork(); \}
        printf("z\textbackslash n");}\\ print? Draw the tree.
  \item A child calls \texttt{execlp} on a program name that does not exist.
        What happens to the code after the call?
  \item Distinguish a zombie from an orphan, and state who eventually reaps
        each.
\end{enumerate}

\subsection*{Solutions}

\begingroup\small
\begin{enumerate}
  \item The child sees \texttt{0}; the parent sees the child's positive pid;
        on failure the parent sees \texttt{-1} and no child exists. After a
        successful \texttt{fork} the two processes run the same code from the
        same point with identical memory, so the only datum that differs is
        the value the call returned.
  \item The \texttt{printf} is after the loop, so we count processes. Four
        sequential forks give $2^4 = 16$ processes, so 16 lines are printed,
        and $16 - 1 = 15$ processes are new children. (Verified: 16 lines.)
  \item First \texttt{fork}: P and C1. C1 sees 0 and takes the \texttt{if}
        branch, forking once: C1 and C2. P takes the \texttt{else} branch and
        forks twice: after the first, P and C3; both then execute the second
        \texttt{fork}, giving P, C3, C4, C5. The processes reaching
        \texttt{printf} are C1, C2, P, C3, C4, C5: 6 lines. (Verified: 6.)
  \item \texttt{execlp} fails and returns \texttt{-1}; the old image is
        intact, so the code after the call --- normally a \texttt{perror} and
        \texttt{exit(1)} --- runs.
  \item A zombie has terminated but its parent has not yet called
        \texttt{wait}; its parent reaps it with \texttt{wait}. An orphan is
        still running (or will terminate) after its parent exited without
        waiting; it is re-parented to \texttt{init}, which reaps it.
\end{enumerate}
\endgroup

\bibliographystyle{plain}
\bibliography{../../Bibliography_base}

\end{document}
